\chapter{KẾT QUẢ VÀ THẢO LUẬN}

\section{Kết quả tổng hợp}

Sáu mô hình --- XGBoost, XLinear, iTransformer, E-STGCN, ST-XLinear và Ensemble --- được đánh giá đồng thời trên 30 cấu hình thí nghiệm, bao gồm năm tầm nhìn dự báo ($T{+}1, T{+}3, T{+}6, T{+}12, T{+}24$ giờ), ba chiến lược chia dữ liệu (block~5, block~7, block~30 ngày) và hai cụm vùng địa lý (North gồm 6 trạm miền Bắc, South gồm 6 trạm miền Nam). Bốn chỉ số MAE, RMSE, $R^2$ và MAPE được sử dụng song song nhằm đánh giá ở nhiều khía cạnh: sai số tuyệt đối trung bình, mức độ nhạy cảm với giá trị cực đoan, tỷ lệ phương sai được giải thích, và sai số tương đối theo phần trăm. Bảng~\ref{tab:block5_north} đến Bảng~\ref{tab:block30_south} trình bày toàn bộ kết quả đo lường.

%% =========================================
%% BLOCK 5
%% =========================================

\begin{table}[H]
\centering
\caption{Kết quả đánh giá trên block~5 --- North}
\label{tab:block5_north}
\small
\begin{tabular}{l|cccc|cccc|cccc}
\toprule
 & \multicolumn{4}{c|}{T+1} & \multicolumn{4}{c|}{T+6} & \multicolumn{4}{c}{T+24} \\
\textbf{Mô hình} & MAE & RMSE & $R^2$ & MAPE & MAE & RMSE & $R^2$ & MAPE & MAE & RMSE & $R^2$ & MAPE \\
\midrule
XGBoost     & \textbf{6,31} & \textbf{12,98} & \textbf{0,83} & \textbf{22,91} & \textbf{10,99} & \textbf{18,70} & \textbf{0,64} & \textbf{40,50} & \textbf{12,67} & \textbf{20,34} & \textbf{0,54} & \textbf{49,25} \\
XLinear     & 6,51 & 13,67 & 0,81 & 26,88 & 14,28 & 23,38 & 0,43 & 63,85 & 17,29 & 26,99 & 0,25 & 81,96 \\
iTransformer & 6,57 & 14,68 & 0,78 & 23,78 & 14,46 & 24,18 & 0,40 & 55,51 & 17,79 & 27,99 & 0,19 & 72,91 \\
E-STGCN     & 9,34 & 17,37 & 0,68 & 32,99 & 14,02 & 23,89 & 0,41 & 52,99 & 17,15 & 27,70 & 0,21 & 69,34 \\
ST-XLinear  & 6,57 & 13,89 & 0,80 & 24,03 & 13,56 & 23,21 & 0,44 & 51,81 & 16,40 & 26,98 & 0,25 & 66,71 \\
Ensemble    & 6,49 & 13,67 & 0,81 & 26,82 & 14,24 & 23,31 & 0,44 & 64,80 & 17,08 & 26,64 & 0,27 & 85,29 \\
\bottomrule
\end{tabular}
\end{table}

\begin{table}[H]
\centering
\caption{Kết quả đánh giá trên block~5 --- South}
\label{tab:block5_south}
\small
\begin{tabular}{l|cccc|cccc|cccc}
\toprule
 & \multicolumn{4}{c|}{T+1} & \multicolumn{4}{c|}{T+6} & \multicolumn{4}{c}{T+24} \\
\textbf{Mô hình} & MAE & RMSE & $R^2$ & MAPE & MAE & RMSE & $R^2$ & MAPE & MAE & RMSE & $R^2$ & MAPE \\
\midrule
XGBoost     & \textbf{4,04} & \textbf{15,74} & \textbf{0,99} & \textbf{21,25} & \textbf{8,18} & \textbf{36,54} & \textbf{0,93} & \textbf{36,76} & \textbf{13,16} & \textbf{71,10} & \textbf{0,73} & \textbf{54,26} \\
XLinear     & 5,47 & 28,63 & 0,96 & 29,65 & 12,46 & 59,43 & 0,83 & 69,99 & 22,67 & 105,66 & 0,45 & 116,00 \\
iTransformer & 5,45 & 32,53 & 0,95 & 24,72 & 14,28 & 71,81 & 0,75 & 58,82 & 21,77 & 102,76 & 0,48 & 104,82 \\
E-STGCN     & 8,79 & 42,31 & 0,91 & 33,35 & 14,17 & 70,86 & 0,75 & 61,53 & 21,74 & 103,06 & 0,48 & 91,53 \\
ST-XLinear  & 5,54 & 30,96 & 0,95 & 24,82 & 13,21 & 66,67 & 0,78 & 61,92 & 22,72 & 103,09 & 0,48 & 83,58 \\
Ensemble    & 5,49 & 28,70 & 0,96 & 29,72 & 12,47 & 59,53 & 0,83 & 69,94 & 22,64 & 105,35 & 0,45 & 116,97 \\
\bottomrule
\end{tabular}
\end{table}

%% =========================================
%% BLOCK 7
%% =========================================

\begin{table}[H]
\centering
\caption{Kết quả đánh giá trên block~7 --- North}
\label{tab:block7_north}
\small
\begin{tabular}{l|cccc|cccc|cccc}
\toprule
 & \multicolumn{4}{c|}{T+1} & \multicolumn{4}{c|}{T+6} & \multicolumn{4}{c}{T+24} \\
\textbf{Mô hình} & MAE & RMSE & $R^2$ & MAPE & MAE & RMSE & $R^2$ & MAPE & MAE & RMSE & $R^2$ & MAPE \\
\midrule
XGBoost     & \textbf{6,17} & \textbf{12,66} & \textbf{0,81} & \textbf{22,26} & \textbf{10,93} & \textbf{19,12} & \textbf{0,61} & \textbf{39,13} & \textbf{12,87} & \textbf{21,24} & \textbf{0,52} & \textbf{51,45} \\
XLinear     & 6,37 & 13,29 & 0,79 & 26,94 & 13,45 & 21,55 & 0,43 & 58,83 & 17,73 & 26,12 & 0,15 & 87,83 \\
iTransformer & 6,45 & 13,99 & 0,76 & 23,94 & 13,92 & 22,62 & 0,38 & 55,64 & 18,50 & 27,50 & 0,06 & 84,61 \\
E-STGCN     & 8,90 & 15,73 & 0,70 & 32,08 & 13,49 & 22,33 & 0,39 & 51,69 & 17,05 & 26,39 & 0,13 & 69,47 \\
ST-XLinear  & 6,39 & 13,27 & 0,79 & 23,81 & 13,02 & 21,67 & 0,43 & 51,32 & 16,51 & 25,63 & 0,18 & 69,95 \\
Ensemble    & 6,35 & 13,28 & 0,79 & 26,84 & 13,40 & 21,50 & 0,44 & 59,89 & 17,26 & 25,72 & 0,18 & 89,22 \\
\bottomrule
\end{tabular}
\end{table}

\begin{table}[H]
\centering
\caption{Kết quả đánh giá trên block~7 --- South}
\label{tab:block7_south}
\small
\begin{tabular}{l|cccc|cccc|cccc}
\toprule
 & \multicolumn{4}{c|}{T+1} & \multicolumn{4}{c|}{T+6} & \multicolumn{4}{c}{T+24} \\
\textbf{Mô hình} & MAE & RMSE & $R^2$ & MAPE & MAE & RMSE & $R^2$ & MAPE & MAE & RMSE & $R^2$ & MAPE \\
\midrule
XGBoost     & \textbf{4,21} & \textbf{14,51} & \textbf{0,99} & \textbf{22,26} & \textbf{8,35} & \textbf{32,96} & \textbf{0,94} & \textbf{35,22} & \textbf{11,42} & \textbf{56,88} & \textbf{0,73} & \textbf{47,67} \\
XLinear     & 5,95 & 32,33 & 0,95 & 30,42 & 13,85 & 66,08 & 0,78 & 70,78 & 26,45 & 126,08 & 0,21 & 120,80 \\
iTransformer & 5,18 & 33,50 & 0,94 & 25,64 & 14,33 & 75,44 & 0,72 & 66,08 & 25,41 & 123,53 & 0,24 & 109,44 \\
E-STGCN     & 7,73 & 37,26 & 0,93 & 32,67 & 13,99 & 73,19 & 0,73 & 63,07 & 24,53 & 116,78 & 0,32 & 97,52 \\
ST-XLinear  & 5,52 & 31,37 & 0,95 & 25,32 & 16,00 & 79,68 & 0,68 & 64,07 & 24,00 & 111,18 & 0,39 & 84,45 \\
Ensemble    & 6,03 & 32,05 & 0,95 & 32,53 & 13,88 & 66,25 & 0,78 & 70,85 & 26,45 & 126,08 & 0,21 & 120,80 \\
\bottomrule
\end{tabular}
\end{table}

%% =========================================
%% BLOCK 30
%% =========================================

\begin{table}[H]
\centering
\caption{Kết quả đánh giá trên block~30 --- North}
\label{tab:block30_north}
\small
\begin{tabular}{l|cccc|cccc|cccc}
\toprule
 & \multicolumn{4}{c|}{T+1} & \multicolumn{4}{c|}{T+6} & \multicolumn{4}{c}{T+24} \\
\textbf{Mô hình} & MAE & RMSE & $R^2$ & MAPE & MAE & RMSE & $R^2$ & MAPE & MAE & RMSE & $R^2$ & MAPE \\
\midrule
XGBoost     & 6,42 & 13,87 & 0,81 & 21,51 & \textbf{10,67} & \textbf{18,77} & \textbf{0,66} & \textbf{37,82} & \textbf{12,78} & \textbf{22,12} & \textbf{0,55} & \textbf{45,03} \\
XLinear     & 5,88 & \textbf{13,38} & \textbf{0,82} & 21,97 & 12,46 & 21,25 & 0,56 & 50,87 & 14,91 & 24,79 & 0,40 & 63,53 \\
iTransformer & 5,96 & 14,25 & 0,80 & 20,04 & 12,07 & 21,38 & 0,55 & 41,26 & 15,05 & 24,82 & 0,40 & 56,91 \\
E-STGCN     & 9,14 & 17,02 & 0,72 & 31,47 & 12,53 & 21,78 & 0,53 & 45,03 & 14,34 & 24,58 & 0,41 & 52,83 \\
ST-XLinear  & \textbf{5,84} & 13,55 & \textbf{0,82} & \textbf{19,49} & 11,81 & 21,86 & 0,53 & 40,72 & 14,37 & 25,42 & 0,37 & 49,51 \\
Ensemble    & 5,88 & 13,39 & \textbf{0,82} & 21,88 & 12,41 & 21,16 & 0,56 & 51,41 & 14,91 & 24,84 & 0,40 & 64,04 \\
\bottomrule
\end{tabular}
\end{table}

\begin{table}[H]
\centering
\caption{Kết quả đánh giá trên block~30 --- South}
\label{tab:block30_south}
\small
\begin{tabular}{l|cccc|cccc|cccc}
\toprule
 & \multicolumn{4}{c|}{T+1} & \multicolumn{4}{c|}{T+6} & \multicolumn{4}{c}{T+24} \\
\textbf{Mô hình} & MAE & RMSE & $R^2$ & MAPE & MAE & RMSE & $R^2$ & MAPE & MAE & RMSE & $R^2$ & MAPE \\
\midrule
XGBoost     & 4,40 & 23,60 & \textbf{0,97} & 21,01 & \textbf{8,38} & \textbf{40,04} & \textbf{0,92} & \textbf{36,60} & 17,68 & 96,61 & 0,55 & \textbf{44,40} \\
XLinear     & \textbf{4,10} & 23,49 & \textbf{0,97} & 24,81 & 11,51 & 54,31 & 0,84 & 55,78 & 16,85 & 89,47 & 0,57 & 93,48 \\
iTransformer & 4,42 & 29,41 & 0,95 & \textbf{20,17} & 10,81 & 63,91 & 0,78 & 50,26 & \textbf{15,33} & 85,26 & \textbf{0,61} & 65,49 \\
E-STGCN     & 8,37 & 47,00 & 0,88 & 31,44 & 11,78 & 63,27 & 0,79 & 51,69 & 15,80 & \textbf{85,18} & \textbf{0,61} & 63,45 \\
ST-XLinear  & 4,33 & \textbf{23,43} & \textbf{0,97} & 20,59 & 11,41 & 54,02 & 0,84 & 51,31 & 19,00 & 91,65 & 0,55 & 65,35 \\
Ensemble    & 4,12 & 23,55 & \textbf{0,97} & 24,91 & 11,53 & 54,42 & 0,84 & 55,90 & 16,81 & 88,96 & 0,58 & 94,65 \\
\bottomrule
\end{tabular}
\end{table}

Từ sáu bảng kết quả, ba xu hướng được quan sát xuyên suốt toàn bộ không gian thí nghiệm.

Xu hướng thứ nhất được quan sát là hiệu suất suy giảm đơn điệu khi tầm nhìn dự báo mở rộng, tuy nhiên tốc độ suy giảm khác nhau đáng kể giữa các kiến trúc. Trên cả sáu bảng, MAE và RMSE tăng liên tục từ $T{+}1$ đến $T{+}24$, đồng thời $R^2$ giảm dần --- tuân theo đặc trưng cơ bản của bài toán dự báo chuỗi thời gian, khi thông tin từ quá khứ bị suy giảm theo khoảng cách tới mốc dự báo. Tuy nhiên, biên độ suy giảm giữa các mô hình dao động lớn: tại block~7 North, chênh lệch MAE giữa $T{+}1$ và $T{+}24$ của XGBoost là 6,70~$\mu g/m^3$ (từ 6,17 lên 12,87), trong khi con số tương ứng của iTransformer là 12,05~$\mu g/m^3$ (từ 6,45 lên 18,50) --- gấp 1,8 lần. Ngược lại, E-STGCN ghi nhận biên độ chỉ 8,15~$\mu g/m^3$ (từ 8,90 lên 17,05), mặc dù giá trị khởi đầu tại $T{+}1$ cao hơn các mô hình khác. Sự khác biệt này phản ánh vai trò của kiến trúc mô hình trong việc duy trì chất lượng dự báo: các kiến trúc có khả năng khai thác thông tin bổ sung (đặc trưng kỹ thuật ở XGBoost, thông tin spatial ở E-STGCN) có xu hướng suy giảm chậm hơn.

Xu hướng thứ hai là tồn tại chênh lệch hiệu suất mang tính hệ thống giữa hai vùng địa lý. So sánh giữa các cặp bảng cùng block (Bảng~\ref{tab:block5_north} với Bảng~\ref{tab:block5_south}; Bảng~\ref{tab:block7_north} với Bảng~\ref{tab:block7_south}; Bảng~\ref{tab:block30_north} với Bảng~\ref{tab:block30_south}), RMSE tại South luôn cao hơn North, với tỷ lệ dao động từ 1,1 lần tại $T{+}1$ đến 2,7--4,5 lần tại $T{+}24$ tùy mô hình. Đồng thời, MAE tại $T{+}1$ South lại thấp hơn North (4,04--8,79 so với 5,84--9,34 trên block~5) --- một hiện tượng được giải thích bởi đặc điểm phân phối: khu vực miền Nam có giá trị nền $PM_{2.5}$ thấp phổ biến (dẫn đến MAE thấp ở dự báo ngắn hạn), nhưng lại xuất hiện các đợt ô nhiễm cực đoan với nồng độ vượt 200~$\mu g/m^3$ (kéo RMSE lên cao). $R^2$ tại South cũng cao hơn North ở $T{+}1$ (đạt 0,88--0,99 so với 0,68--0,83), phản ánh phương sai tổng lớn của dữ liệu phía Nam theo công thức $R^2 = 1 - SS_{res}/SS_{tot}$.

Xu hướng thứ ba là thứ tự xếp hạng giữa các mô hình thay đổi tùy theo chiến lược chia dữ liệu. Tại block~5 và block~7, XGBoost ghi nhận MAE thấp nhất trên hầu hết cấu hình, đặc biệt ở horizon dài. Tuy nhiên, tại block~30 North $T{+}1$, XLinear (MAE~$= 5{,}88$), ST-XLinear (MAE~$= 5{,}84$) và Ensemble (MAE~$= 5{,}88$) đều ghi nhận MAE thấp hơn XGBoost (MAE~$= 6{,}42$). Tương tự, tại block~30 South $T{+}24$, iTransformer ($R^2 = 0{,}61$) và E-STGCN ($R^2 = 0{,}61$) đạt $R^2$ cao hơn XGBoost ($R^2 = 0{,}55$). Hiện tượng đảo thứ hạng này cho thấy hiệu suất tương đối giữa các mô hình phụ thuộc vào lượng dữ liệu huấn luyện liên tục: block~30 cung cấp khoảng 720 bước thời gian liên tục cho mỗi khối, gấp 6 lần block~5 (khoảng 120 bước), và các kiến trúc deep learning (attention, graph convolution) --- vốn cần nhiều mẫu để ước lượng tham số --- hưởng lợi rõ rệt từ sự gia tăng này. Mối quan hệ giữa kiến trúc mô hình và chiến lược chia dữ liệu được phân tích chi tiết trong Mục~\ref{sec:block_comparison}.

\section{Phân tích hiệu suất theo tầm nhìn dự báo}

Mục này phân tích đặc trưng suy giảm hiệu suất của từng mô hình khi tầm nhìn dự báo mở rộng từ $T{+}1$ đến $T{+}24$ giờ, thông qua từng chỉ số đánh giá. Toàn bộ biểu đồ sử dụng định dạng $2 \times 3$ ô con (hai hàng tương ứng miền Bắc và miền Nam, ba cột tương ứng block~5, block~7 và block~30) với cùng thang đo trục tung, cho phép so sánh trực quan giữa các cấu hình mà không bị ảnh hưởng bởi sự khác biệt tỷ lệ.

\subsection{Sai số tuyệt đối trung bình (MAE)}

Các đường MAE theo tầm nhìn dự báo phân kỳ rõ rệt từ $T{+}6$ trở đi, trong khi hội tụ chặt tại $T{+}1$. Hình~\ref{fig:mae_vs_horizon} minh họa hiện tượng này.

\begin{figure}[H]
\centering
\includegraphics[width=\textwidth]{metric_mae_vs_horizon.png}
\caption{MAE theo tầm nhìn dự báo --- 3 block $\times$ 2 vùng (cùng thang đo)}
\label{fig:mae_vs_horizon}
\end{figure}

Tại $T{+}1$, sáu mô hình hội tụ trong khoảng 4,04--9,34~$\mu g/m^3$, với chênh lệch giữa mô hình có MAE thấp nhất và cao nhất chỉ khoảng 3--5~$\mu g/m^3$. Sự hội tụ này phản ánh bản chất của dự báo 1 giờ phía trước: quán tính chuỗi thời gian cung cấp phần lớn thông tin cần thiết, và mọi kiến trúc đều khai thác được nguồn tín hiệu này. Từ $T{+}6$ trở đi, quán tính suy yếu và các đường bắt đầu tách biệt: tại block~7 miền Bắc $T{+}6$, khoảng cách giữa XGBoost (10,93) và iTransformer (13,92) đã đạt 2,99~$\mu g/m^3$, và mở rộng lên 5,63~$\mu g/m^3$ tại $T{+}24$ (12,87 so với 18,50). Khoảng cách mở rộng này cho thấy ở tầm nhìn dài, thông tin từ hệ thống xây dựng đặc trưng thủ công (độ trễ, trung bình trượt, thời tiết tương lai) mà XGBoost sử dụng ngày càng có giá trị dự báo cao hơn so với các biểu diễn tự học của kiến trúc học sâu.

Tại miền Nam, MAE tuyệt đối ở $T{+}1$ thấp hơn miền Bắc do phân phối $PM_{2.5}$ phía Nam có giá trị nền thấp phổ biến, tuy nhiên tốc độ tăng tuyệt đối lớn hơn khi tầm nhìn mở rộng. Cụ thể tại block~7, XGBoost tăng từ 4,21 lên 11,42 ($\Delta = 7{,}21$~$\mu g/m^3$), trong khi XLinear tăng từ 5,95 lên 26,45 ($\Delta = 20{,}50$~$\mu g/m^3$) --- gấp gần ba lần. Khoảng cách $\Delta$ lớn này phản ánh tính chất dữ liệu miền Nam: các đợt ô nhiễm cực đoan ngắn hạn tạo ra biến thiên phi tuyến, và các mô hình không có thông tin thời tiết tương lai hoặc đặc trưng không gian gặp khó khăn trong việc dự báo chính xác các sự kiện này.

E-STGCN thể hiện đặc trưng riêng biệt đáng chú ý: mặc dù khởi đầu với MAE cao nhất ở $T{+}1$ (8,90--9,34 tại miền Bắc, 7,73--8,79 tại miền Nam), mô hình này có tốc độ tăng MAE chậm nhất trong nhóm học sâu. Tại block~30 miền Bắc, MAE $T{+}24$ của E-STGCN (14,34) thấp hơn XLinear (14,91), iTransformer (15,05) và ST-XLinear (14,37). Hiện tượng này gợi ý rằng kiến trúc tích chập đồ thị kết hợp LSTM cho phép E-STGCN khai thác tương quan không gian giữa các trạm, bù đắp cho sự thiếu hụt thông tin thời gian ở tầm nhìn xa, nhưng đồng thời cũng gây ra sai lệch hệ thống ở $T{+}1$ do cơ chế lan truyền đồ thị làm trung bình hóa tín hiệu cục bộ.

\subsection{Sai số toàn phương trung bình (RMSE)}

RMSE khuếch đại sự khác biệt giữa hai vùng lên mức đáng kể so với MAE, do phép bình phương trong công thức RMSE phạt nặng các sai lệch lớn. Hình~\ref{fig:rmse_vs_horizon} trình bày kết quả.

\begin{figure}[H]
\centering
\includegraphics[width=\textwidth]{metric_rmse_vs_horizon.png}
\caption{RMSE theo tầm nhìn dự báo --- 3 block $\times$ 2 vùng (cùng thang đo)}
\label{fig:rmse_vs_horizon}
\end{figure}

Tại block~7 $T{+}24$, tỷ lệ RMSE giữa miền Nam và miền Bắc đạt 2,7 lần cho XGBoost (56,88/21,24) và 4,5 lần cho iTransformer (123,53/27,50). Trong khi đó, tỷ lệ MAE tương ứng chỉ khoảng 1,4--1,9 lần. Sự chênh lệch giữa hai tỷ lệ này cho thấy vấn đề ở miền Nam không chỉ là sai số trung bình lớn hơn mà còn là tần suất xuất hiện các sai lệch cực lớn cao hơn --- tức mô hình tích lũy sai số khi gặp các đợt ô nhiễm với nồng độ vượt xa phân phối huấn luyện. XGBoost kiểm soát hiện tượng này nhờ cơ chế tăng cường gradient tuần tự: mỗi cây quyết định mới được huấn luyện tập trung vào phần dư lớn nhất từ vòng trước, tạo ra khả năng tự hiệu chỉnh đối với giá trị cực đoan.

So sánh giữa ba block, RMSE tại miền Nam giảm rõ rệt khi tăng kích thước block: XLinear $T{+}24$ giảm từ 105,66 (block~5) xuống 126,08 (block~7) rồi xuống 89,47 (block~30). Xu hướng giảm ở block~30 cho thấy tập huấn luyện liên tục dài hơn giúp mô hình học sâu tiếp xúc với nhiều đợt ô nhiễm cực đoan trong quá trình huấn luyện, từ đó cải thiện khả năng xử lý các sự kiện tương tự trên tập kiểm thử. Tuy nhiên, giá trị RMSE tuyệt đối của nhóm học sâu tại miền Nam vẫn cao hơn XGBoost 2--3 lần ngay cả ở block~30 (89,47 so với 96,61 cho XGBoost tại block~30 miền Nam $T{+}24$), cho thấy khoảng cách chưa được thu hẹp hoàn toàn.

\subsection{Hệ số xác định $R^2$}

Mức độ giải thích phương sai dữ liệu suy giảm nhanh ở tầm nhìn dài, đặc biệt trên block~5 và block~7, trong khi block~30 duy trì $R^2$ ở mức cao hơn có ý nghĩa. Hình~\ref{fig:r2_heatmap} thể hiện $R^2$ dưới dạng bản đồ nhiệt, cho phép nhận diện nhanh vùng hiệu suất thấp (ô màu đỏ) và vùng hiệu suất cao (ô màu xanh).

\begin{figure}[H]
\centering
\includegraphics[width=\textwidth]{r2_heatmap_all.png}
\caption{Bản đồ nhiệt $R^2$ --- tất cả block $\times$ vùng (thang đo $0$--$1$)}
\label{fig:r2_heatmap}
\end{figure}

Bản đồ nhiệt bộc lộ chiều hướng suy giảm theo chiều ngang rõ rệt ở mọi cấu hình. Tại block~7 miền Bắc, $R^2$ của iTransformer giảm từ 0,76 ($T{+}1$) xuống 0,06 ($T{+}24$), tức mô hình chỉ còn giải thích được 6\% phương sai dữ liệu --- gần tương đương với dự báo bằng giá trị trung bình lịch sử. Trong cùng cấu hình, XGBoost giữ $R^2 = 0{,}52$, vẫn giải thích được hơn một nửa biến thiên. Sự chênh lệch 0,46 giữa hai mô hình tại cùng một cấu hình phản ánh khác biệt cơ bản trong kiến trúc: cơ chế attention đảo ngược của iTransformer phụ thuộc vào khả năng tập trung trọng số vào các bước thời gian có tính dự báo cao, nhưng khi mốc dự báo xa dần khỏi cửa sổ quan sát, cơ chế này mất dần tính chọn lọc và phân tán trọng số đều trên toàn bộ chuỗi.

Tại block~30, sự suy giảm được kiềm chế đáng kể. $R^2$ tại $T{+}24$ miền Bắc của iTransformer tăng từ 0,06 (block~7) lên 0,40 (block~30), tức tăng gấp 6,7 lần; E-STGCN tăng từ 0,13 lên 0,41, tức tăng gấp 3,2 lần. Bước nhảy này chứng minh rằng vấn đề chính của hai kiến trúc này trên block~5 và block~7 không phải do bản chất kiến trúc kém mà do thiếu dữ liệu huấn luyện liên tục. Cơ chế attention cần đủ mẫu để học phân phối trọng số phù hợp, và tích chập đồ thị cần đủ ngữ cảnh lịch sử để ước lượng ma trận trọng số giữa các trạm.

Một hiện tượng cần lưu ý khi đọc bản đồ nhiệt là $R^2$ tại miền Nam ở $T{+}1$ luôn cao hơn miền Bắc (0,88--0,99 so với 0,68--0,83), mặc dù RMSE tại miền Nam cao hơn. Hiện tượng này không phản ánh chất lượng dự báo tốt hơn mà do đặc tính phân phối: phía Nam có phương sai tổng ($SS_{tot}$) lớn --- dải giá trị $PM_{2.5}$ trải rộng từ gần 0 đến trên 200~$\mu g/m^3$ --- nên theo công thức $R^2 = 1 - SS_{res}/SS_{tot}$, dù $SS_{res}$ lớn, tỷ lệ $SS_{res}/SS_{tot}$ vẫn nhỏ. Do đó, khi so sánh chất lượng dự báo tuyệt đối giữa hai vùng, MAE và RMSE là chỉ số phù hợp hơn $R^2$.

\subsection{Sai số phần trăm tuyệt đối trung bình (MAPE)}

Toàn bộ mô hình ghi nhận MAPE ở mức 19--121\% trên mọi cấu hình, một dải giá trị cần được giải thích trước khi diễn giải kết quả. Hình~\ref{fig:mape_vs_horizon} trình bày kết quả.

\begin{figure}[H]
\centering
\includegraphics[width=\textwidth]{metric_mape_vs_horizon.png}
\caption{MAPE theo tầm nhìn dự báo --- 3 block $\times$ 2 vùng (cùng thang đo)}
\label{fig:mape_vs_horizon}
\end{figure}

MAPE cao không phản ánh chất lượng dự báo kém mà phản ánh đặc tính cố hữu của phân phối $PM_{2.5}$. Nồng độ $PM_{2.5}$ tại nhiều thời điểm có giá trị rất thấp (dưới 10~$\mu g/m^3$, phổ biến vào ban đêm và mùa hè), khiến sai lệch tuyệt đối chỉ 2--3~$\mu g/m^3$ cũng tạo ra MAPE hàng chục phần trăm. Ví dụ, với giá trị thực $5~\mu g/m^3$ và dự báo $8~\mu g/m^3$, sai lệch tuyệt đối chỉ 3~$\mu g/m^3$ nhưng MAPE đạt 60\%. Vì lý do này, MAPE trong bài toán dự báo $PM_{2.5}$ cần được xem xét kết hợp với MAE và RMSE, không nên đánh giá độc lập.

Dù giá trị tuyệt đối cao, thứ tự xếp hạng MAPE vẫn nhất quán với các chỉ số khác: XGBoost ghi nhận MAPE thấp nhất (21--54\%), tiếp theo là ST-XLinear và iTransformer. Một ngoại lệ đáng chú ý xuất hiện tại block~30 $T{+}1$: iTransformer đạt MAPE thấp nhất (20,04\% North, 20,17\% South), trong khi XGBoost đạt 21,51\% và 21,01\%. Sự đảo ngược này cho thấy cơ chế attention có thể tập trung chính xác vào bước thời gian liền kề --- phát huy thế mạnh ở dự báo ngắn hạn --- khi tập huấn luyện đủ lớn để ước lượng phân phối trọng số attention. Ensemble ghi nhận MAPE cao nhất ở nhiều cấu hình (lên đến 120,80\% tại block~7 South $T{+}24$), do phương pháp trung bình trọng số kế thừa và phóng đại sai lệch phần trăm từ các mô hình thành phần khi mẫu số (giá trị thực) nhỏ.

\section{Đánh giá tính ổn định}

Tính ổn định của mô hình được đánh giá qua ba góc nhìn bổ sung nhau: khoảng cách hiệu suất giữa dự báo ngắn hạn và dài hạn, phân bố sai số qua các cấu hình, và mức độ cân bằng đa metric. Ba phân tích này giúp trả lời câu hỏi: mô hình nào duy trì chất lượng dự báo ổn định nhất khi điều kiện thí nghiệm thay đổi?

\subsection{Khoảng cách hiệu suất ngắn hạn -- dài hạn}

Khoảng cách MAE giữa $T{+}1$ và $T{+}24$ (ký hiệu $\Delta$MAE $=$ MAE$_{T+24} -$ MAE$_{T+1}$) đo lường mức suy giảm hiệu suất khi mở rộng tầm nhìn. Mô hình có $\Delta$ nhỏ duy trì chất lượng dự báo ổn định xuyên suốt phổ tầm nhìn. Hình~\ref{fig:mae_gap} trình bày so sánh trực tiếp MAE tại hai mốc, và Bảng~\ref{tab:delta_mae} tổng hợp $\Delta$ trên toàn bộ cấu hình.

\begin{figure}[H]
\centering
\includegraphics[width=\textwidth]{mae_t1_vs_t24_gap.png}
\caption{Khoảng cách MAE giữa $T{+}1$ và $T{+}24$ --- cùng thang đo}
\label{fig:mae_gap}
\end{figure}

\begin{table}[H]
\centering
\caption{$\Delta$MAE ($=$ MAE$_{T+24} -$ MAE$_{T+1}$) trên tất cả cấu hình ($\mu g/m^3$)}
\label{tab:delta_mae}
\small
\begin{tabular}{l|cc|cc|cc}
\toprule
 & \multicolumn{2}{c|}{Block 5} & \multicolumn{2}{c|}{Block 7} & \multicolumn{2}{c}{Block 30} \\
\textbf{Mô hình} & Bắc & Nam & Bắc & Nam & Bắc & Nam \\
\midrule
XGBoost     & \textbf{6,36} & \textbf{9,12} & \textbf{6,70} & \textbf{7,21} & 6,36 & 13,28 \\
XLinear     & 10,78 & 17,20 & 11,36 & 20,50 & 9,03 & 12,75 \\
iTransformer & 11,22 & 16,32 & 12,05 & 20,23 & 9,09 & 10,91 \\
E-STGCN     & 7,81 & 12,95 & 8,15 & 16,80 & \textbf{5,20} & \textbf{7,43} \\
ST-XLinear  & 9,83 & 17,18 & 10,12 & 18,48 & 8,53 & 14,67 \\
Ensemble    & 10,59 & 17,15 & 10,91 & 20,42 & 9,03 & 12,69 \\
\bottomrule
\end{tabular}
\end{table}

Bảng~\ref{tab:delta_mae} cho thấy mỗi kiến trúc có đặc trưng $\Delta$ riêng, phụ thuộc vào cả chiến lược chia dữ liệu và vùng địa lý. Trên block~5 và block~7, XGBoost ghi nhận $\Delta$ thấp nhất trên cả bốn cấu hình: 6,36--6,70 tại North và 7,21--9,12 tại South. Giá trị $\Delta$ tại North dao động hẹp (chênh lệch chỉ 0,34 giữa block~5 và block~7) cho thấy XGBoost ít nhạy cảm với chiến lược chia dữ liệu ở vùng có phân phối ổn định.

Tại block~30, thứ tự xếp hạng thay đổi: E-STGCN ghi nhận $\Delta$ thấp nhất (5,20 miền Bắc, 7,43 miền Nam), thấp hơn XGBoost (6,36 và 13,28). Nguyên nhân là tập huấn luyện liên tục 30 ngày cung cấp đủ mẫu để tích chập đồ thị học được ma trận kề phản ánh tương quan giữa các trạm --- lợi thế không gian này bù đắp sự thiếu hụt đặc trưng thủ công, và tác dụng bù đắp thể hiện mạnh nhất ở tầm nhìn dài nơi thông tin thời gian đã suy yếu.

$\Delta$ tại South luôn lớn hơn North trên mọi mô hình, với tỷ lệ $\Delta_{\text{South}}/\Delta_{\text{North}}$ dao động 1,1--2,1. Tỷ lệ này lớn nhất ở block~7 (1,8--2,0) và nhỏ nhất ở block~30 (1,2--1,7), cho thấy tập huấn luyện lớn giúp thu hẹp khoảng cách ổn định giữa hai vùng. Hiện tượng này phản ánh việc miền Nam có nhiều sự kiện cực đoan hơn --- vốn là nguồn gây mất ổn định chính --- và block~30 cho phép mô hình tiếp xúc đủ với các sự kiện này trong quá trình huấn luyện.

\subsection{Phân bố RMSE qua các cấu hình}

Phân bố RMSE dưới dạng biểu đồ hộp trên lưới $3 \times 5$ (block $\times$ tầm nhìn) bổ sung cho phân tích $\Delta$MAE bằng cách thể hiện cả trung vị lẫn độ trải (khoảng tứ phân vị), cho phép đánh giá đồng thời giá trị trung tâm và mức độ phân tán. Hình~\ref{fig:boxplot} trình bày kết quả với cùng trục tung trên 15 ô con.

\begin{figure}[H]
\centering
\includegraphics[width=\textwidth]{rmse_boxplot_all_blocks.png}
\caption{Phân bố RMSE --- 3 block $\times$ 5 horizon (cùng thang đo)}
\label{fig:boxplot}
\end{figure}

Biểu đồ bộc lộ hai chiều biến thiên đồng thời. Theo chiều tầm nhìn (từ trái sang phải trong mỗi hàng), trung vị và độ trải của mọi hộp đều tăng, phản ánh sự suy giảm hiệu suất kèm theo gia tăng bất ổn. Tuy nhiên, tốc độ mở quạt (tức tỷ lệ giữa độ trải ở $T{+}24$ và độ trải ở $T{+}1$) khác nhau giữa các mô hình: XGBoost có hình quạt mở chậm và đều, cho thấy sai số tăng nhưng vẫn kiểm soát được; XLinear và iTransformer mở nhanh hơn, phản ánh tính nhạy cảm cao với dữ liệu kiểm thử ở tầm nhìn dài.

Theo chiều block (từ hàng trên xuống dưới), biểu đồ hộp của nhóm học sâu co lại đáng kể từ block~5 xuống block~30, đặc biệt ở $T{+}12$ và $T{+}24$: độ trải giảm, trung vị giảm, và khoảng cách giữa các hộp trở nên hẹp hơn. XGBoost có hộp kích thước ổn định qua ba hàng, cho thấy mô hình này ít chịu ảnh hưởng bởi chiến lược chia dữ liệu.

E-STGCN có hộp với độ trải nhỏ nhưng trung vị cao ở $T{+}1$, rồi trung vị giảm tương đối so với các mô hình khác ở $T{+}24$. Hình dạng này nhất quán với phân tích MAE ở Mục 3.2.1: tích chập đồ thị gây sai lệch hệ thống ở ngắn hạn (do lan truyền đồ thị làm trung bình hóa) nhưng duy trì ổn định ở dài hạn (nhờ thông tin không gian bù đắp).

\subsection{Đánh giá đa metric đồng thời}

Biểu đồ radar đa chỉ số tại $T{+}1$ cung cấp góc nhìn tổng hợp về năng lực dự báo trên nhiều tiêu chí đồng thời. Hình~\ref{fig:radar} trình bày biểu đồ radar trên toàn bộ $2 \times 3$ cấu hình, sử dụng chuẩn hóa min-max trong nhóm (block, vùng) để đưa bốn chỉ số về thang $[0, 1]$ --- giá trị 1 tương ứng mô hình ghi nhận kết quả tốt nhất trên chỉ số đó. Diện tích đa giác phản ánh mức độ cân bằng: mô hình có đa giác lớn và đều đặn trên bốn trục thể hiện năng lực toàn diện, trong khi đa giác méo cho thấy mô hình mạnh trên một số tiêu chí nhưng yếu trên các tiêu chí khác.

\begin{figure}[H]
\centering
\includegraphics[width=\textwidth]{radar_t1_all.png}
\caption{Biểu đồ radar đa chỉ số tại $T{+}1$ --- chuẩn hóa min-max, tất cả cấu hình}
\label{fig:radar}
\end{figure}

Tại block~5 và block~7, XGBoost có đa giác phủ gần đỉnh trên cả bốn trục, phản ánh giá trị gần 1 đồng thời trên MAE, RMSE, $R^2$ và MAPE. Các mô hình XLinear, iTransformer và ST-XLinear có đa giác nhỏ hơn nhưng tương đối đều, cho thấy hiệu suất thấp hơn trên mọi chỉ số nhưng không bị lệch nghiêm trọng về phía chỉ số nào. Ensemble có hình dạng gần trùng với XLinear do cơ chế trung bình trọng số phụ thuộc mạnh vào mô hình thành phần có trọng số cao nhất.

E-STGCN thể hiện hình dạng bất đối xứng nhất quán trên mọi cấu hình: đỉnh trên trục $R^2$ tương đương hoặc gần bằng các mô hình học sâu khác, nhưng đỉnh trên trục MAE và RMSE thấp rõ rệt. Hình dạng này phản ánh đặc trưng kiến trúc đã phân tích ở Mục 3.2.1: tích chập đồ thị nắm bắt tốt xu hướng biến động (thể hiện qua $R^2$) nhưng gây sai lệch hệ thống làm tăng sai số tuyệt đối (thể hiện qua MAE/RMSE).

Tại block~30 miền Bắc, đa giác của các mô hình hội tụ rõ rệt: XLinear, ST-XLinear và Ensemble có diện tích gần bằng XGBoost. Sự hội tụ này cho thấy khi dữ liệu huấn luyện đủ lớn, khoảng cách hiệu suất giữa các kiến trúc thu hẹp ở dự báo ngắn hạn. Tại block~30 miền Nam, sự hội tụ kém hơn do phân phối dữ liệu phức tạp hơn, nhưng iTransformer thu hẹp khoảng cách so với XGBoost trên trục MAPE --- nhất quán với phân tích ở Mục 3.2.4.

\section{Ảnh hưởng của vùng miền}

Chênh lệch hiệu suất giữa miền Bắc và miền Nam mang tính hệ thống và không phụ thuộc vào kiến trúc mô hình. Hình~\ref{fig:north_south_gap} trình bày khoảng cách MAE và RMSE giữa hai vùng trên toàn bộ cấu hình.

\begin{figure}[H]
\centering
\includegraphics[width=\textwidth]{north_south_gap.png}
\caption{Chênh lệch hiệu suất giữa miền Bắc và miền Nam --- tất cả cấu hình}
\label{fig:north_south_gap}
\end{figure}

Biểu đồ cho thấy sáu mô hình đều ghi nhận mẫu hình chênh lệch tương tự: MAE tại miền Bắc cao hơn miền Nam ở $T{+}1$ nhưng thấp hơn ở $T{+}24$, trong khi RMSE tại miền Nam luôn cao hơn ở mọi tầm nhìn. Sự đồng nhất của mẫu hình này trên cả sáu kiến trúc cho thấy nguyên nhân nằm ở đặc điểm dữ liệu chứ không phải ở thiết kế mô hình. Có thể nhận diện ba nguyên nhân chính từ dữ liệu.

Nguyên nhân thứ nhất là sự khác biệt về phân phối nồng độ $PM_{2.5}$. Khu vực miền Bắc có nồng độ nền cao và ổn định hơn (trung bình khoảng 30--60~$\mu g/m^3$), dẫn đến MAE tuyệt đối cao hơn tại $T{+}1$ vì sai lệch được tính trên giá trị nền lớn hơn. Ngược lại, miền Nam có nồng độ nền thấp phổ biến (dưới 20~$\mu g/m^3$) xen lẫn các đợt ô nhiễm cực đoan vượt 200~$\mu g/m^3$, tạo ra phân phối lệch phải với đuôi dài. Phân phối đuôi dài này giải thích hiện tượng MAE thấp nhưng RMSE cao ở miền Nam: phần lớn các điểm dữ liệu có sai số nhỏ (dẫn đến MAE thấp), nhưng một số ít điểm cực đoan có sai lệch rất lớn bị phóng đại bởi phép bình phương trong RMSE.

Nguyên nhân thứ hai liên quan đến đặc trưng khí tượng. Miền Bắc chịu ảnh hưởng rõ rệt của gió mùa Đông Bắc với chu kỳ mùa ổn định, tạo ra tính chu kỳ dễ nắm bắt cho mô hình. Miền Nam có khí hậu nhiệt đới gió mùa với sự thay đổi nhanh giữa mùa khô và mùa mưa, các đợt nghịch nhiệt cục bộ, và hiện tượng đốt rơm rạ theo mùa vụ --- tất cả tạo ra biến thiên phi tuyến khó dự báo ở tầm nhìn dài. Khi tầm nhìn mở rộng từ $T{+}1$ đến $T{+}24$, quán tính chuỗi thời gian suy yếu và các sự kiện khí tượng phi tuyến chiếm ưu thế, khiến sai số tại miền Nam tăng nhanh hơn.

Nguyên nhân thứ ba là mật độ không gian của mạng lưới trạm quan trắc. Sáu trạm miền Bắc tập trung chủ yếu tại vùng đồng bằng sông Hồng với điều kiện địa hình tương đồng, dẫn đến tương quan không gian cao giữa các trạm. Sáu trạm miền Nam phân bố từ đồng bằng đến ven biển, với sự đa dạng về vi khí hậu lớn hơn, làm giảm tương quan không gian. Các mô hình có thành phần không gian (E-STGCN, ST-XLinear) khai thác tương quan này để cải thiện dự báo, nên khoảng cách hiệu suất giữa hai vùng ở nhóm mô hình không gian nhỏ hơn so với nhóm chỉ sử dụng thông tin thời gian (XLinear, iTransformer). Cụ thể tại block~7 $T{+}24$, tỷ lệ $\Delta$MAE(Nam)/$\Delta$MAE(Bắc) của E-STGCN là 2,06, trong khi tỷ lệ tương ứng của XLinear là 1,81 và iTransformer là 1,68.

\section{Ảnh hưởng của chiến lược chia dữ liệu}
\label{sec:block_comparison}

Chiến lược chia dữ liệu theo khối liên tục (block split) ảnh hưởng khác nhau đến từng kiến trúc mô hình. Hình~\ref{fig:block_compare} trình bày so sánh bốn chỉ số giữa ba block trên cùng hệ trục.

\begin{figure}[H]
\centering
\includegraphics[width=\textwidth]{block_compare_mae.png}
\caption{So sánh MAE giữa ba chiến lược chia dữ liệu --- tất cả mô hình}
\label{fig:block_compare}
\end{figure}

\begin{figure}[H]
\centering
\includegraphics[width=\textwidth]{block_compare_rmse.png}
\caption{So sánh RMSE giữa ba chiến lược chia dữ liệu --- tất cả mô hình}
\label{fig:block_compare_rmse}
\end{figure}

\begin{figure}[H]
\centering
\includegraphics[width=\textwidth]{block_compare_r2.png}
\caption{So sánh $R^2$ giữa ba chiến lược chia dữ liệu --- tất cả mô hình}
\label{fig:block_compare_r2}
\end{figure}

\begin{figure}[H]
\centering
\includegraphics[width=\textwidth]{block_compare_mape.png}
\caption{So sánh MAPE giữa ba chiến lược chia dữ liệu --- tất cả mô hình}
\label{fig:block_compare_mape}
\end{figure}

Từ bốn biểu đồ so sánh, nhóm mô hình có thể phân thành hai loại theo mức độ phụ thuộc vào chiến lược chia dữ liệu.

Nhóm thứ nhất gồm XGBoost, có hiệu suất ổn định qua ba block. MAE tại miền Bắc $T{+}24$ dao động trong khoảng 12,67--12,87 qua ba block (biên độ chỉ 0,20~$\mu g/m^3$), và $R^2$ dao động 0,52--0,55 (biên độ 0,03). Tính ổn định này xuất phát từ đặc trưng kiến trúc: XGBoost sử dụng hệ thống đặc trưng thủ công gồm 134 biến được tính toán trước (độ trễ, trung bình trượt, đặc trưng thời gian, thời tiết tương lai), nên không phụ thuộc vào độ dài chuỗi liên tục để học biểu diễn. Mỗi cây quyết định chỉ cần phân chia không gian đặc trưng --- một thao tác không yêu cầu tính liên tục dữ liệu.

Nhóm thứ hai gồm iTransformer, E-STGCN, XLinear, ST-XLinear và Ensemble, có hiệu suất cải thiện đáng kể khi chuyển từ block~5 sang block~30. Mức cải thiện khác nhau tùy kiến trúc:

iTransformer ghi nhận bước nhảy lớn nhất: $R^2$ tại miền Bắc $T{+}24$ tăng từ 0,19 (block~5) lên 0,06 (block~7) rồi lên 0,40 (block~30). Mức tăng 0,34 từ block~7 lên block~30 cho thấy cơ chế attention cần một ngưỡng dữ liệu tối thiểu để hoạt động hiệu quả. Với block~5 (khoảng 120 bước thời gian liên tục) và block~7 (khoảng 168 bước), số lượng mẫu chưa đủ để ước lượng ma trận trọng số attention cho chuỗi dài; block~30 (khoảng 720 bước) vượt qua ngưỡng này. Hiện tượng tương tự được quan sát ở MAE: từ 17,79 (block~5) giảm xuống 15,05 (block~30), tức giảm 15,4\%.

E-STGCN cải thiện đều đặn hơn iTransformer, với mức tăng $R^2$ phân bố tương đối đều giữa ba bước nhảy block. Tại miền Bắc $T{+}24$, $R^2$ tăng từ 0,21 (block~5) lên 0,13 (block~7) rồi lên 0,41 (block~30). Đặc biệt, E-STGCN là mô hình duy nhất ghi nhận $R^2$ tại block~30 cao hơn tất cả các mô hình học sâu khác ở $T{+}24$ miền Bắc (0,41 so với 0,40 của iTransformer, 0,40 của XLinear, 0,37 của ST-XLinear). Điều này cho thấy kiến trúc tích chập đồ thị phát huy tối đa khi có đủ dữ liệu liên tục để ước lượng ma trận kề phản ánh quan hệ không gian giữa các trạm.

XLinear và ST-XLinear cải thiện ở mức trung bình. $R^2$ của XLinear tại miền Bắc $T{+}24$ tăng từ 0,25 (block~5) lên 0,40 (block~30), trong khi ST-XLinear tăng từ 0,25 lên 0,37. Mức cải thiện thấp hơn so với iTransformer phản ánh kiến trúc tuyến tính có ít tham số cần ước lượng hơn, nên hưởng lợi ít hơn từ dữ liệu bổ sung.

Ensemble không cải thiện đáng kể so với mô hình thành phần có hiệu suất cao nhất. Tại block~30 miền Bắc $T{+}24$, $R^2$ của Ensemble (0,40) bằng đúng $R^2$ của XLinear (0,40), cho thấy cơ chế trung bình trọng số tối ưu trên tập xác thực đã gán trọng số gần như toàn bộ cho XLinear. Hiện tượng này xuất hiện nhất quán: Ensemble gần trùng với mô hình thành phần tốt nhất trên tập xác thực thay vì kết hợp bổ sung thông tin từ nhiều nguồn. Nguyên nhân là năm mô hình thành phần chia sẻ cùng nguồn dữ liệu đầu vào và có tương quan sai số cao, khiến phương pháp kết hợp tuyến tính không tạo ra giá trị bổ sung.

Tóm lại, chiến lược chia dữ liệu tác động mạnh nhất đến nhóm mô hình học sâu có nhiều tham số (iTransformer, E-STGCN), tác động trung bình đến nhóm tuyến tính (XLinear, ST-XLinear), và gần như không tác động đến XGBoost. Quan sát này có ý nghĩa thực tiễn quan trọng: khi triển khai hệ thống dự báo chất lượng không khí với dữ liệu hạn chế, XGBoost là lựa chọn ổn định; khi có đủ dữ liệu liên tục (trên 30 ngày mỗi khối), các kiến trúc học sâu trở nên cạnh tranh và trong một số cấu hình còn vượt qua XGBoost.

\section{Tốc độ suy giảm hiệu suất}

Để lượng hóa tốc độ suy giảm hiệu suất, hệ số góc (slope) tuyến tính của MAE theo tầm nhìn dự báo được ước lượng bằng hồi quy tuyến tính đơn giản trên năm điểm ($T{+}1, T{+}3, T{+}6, T{+}12, T{+}24$). Hệ số góc dương cho thấy MAE tăng khi tầm nhìn mở rộng; giá trị slope càng lớn thì mô hình suy giảm càng nhanh. Hình~\ref{fig:degradation_slope} trình bày so sánh trực quan và Bảng~\ref{tab:slope} tổng hợp giá trị slope trên toàn bộ cấu hình.

\begin{figure}[H]
\centering
\includegraphics[width=\textwidth]{degradation_slope.png}
\caption{Hệ số góc suy giảm MAE theo tầm nhìn --- tất cả cấu hình}
\label{fig:degradation_slope}
\end{figure}

\begin{table}[H]
\centering
\caption{Hệ số góc suy giảm MAE (slope, $\mu g/m^3$ mỗi giờ tầm nhìn)}
\label{tab:slope}
\small
\begin{tabular}{l|cc|cc|cc}
\toprule
 & \multicolumn{2}{c|}{Block 5} & \multicolumn{2}{c|}{Block 7} & \multicolumn{2}{c}{Block 30} \\
\textbf{Mô hình} & Bắc & Nam & Bắc & Nam & Bắc & Nam \\
\midrule
XGBoost     & \textbf{0,26} & \textbf{0,37} & \textbf{0,27} & \textbf{0,29} & 0,26 & 0,54 \\
XLinear     & 0,44 & 0,70 & 0,46 & 0,84 & 0,37 & 0,52 \\
iTransformer & 0,46 & 0,66 & 0,49 & 0,82 & 0,37 & 0,44 \\
E-STGCN     & 0,32 & 0,53 & 0,33 & 0,68 & \textbf{0,21} & \textbf{0,30} \\
ST-XLinear  & 0,40 & 0,70 & 0,41 & 0,75 & 0,35 & 0,60 \\
Ensemble    & 0,43 & 0,70 & 0,44 & 0,83 & 0,37 & 0,52 \\
\bottomrule
\end{tabular}
\end{table}

Bảng~\ref{tab:slope} cho thấy sáu mô hình phân thành ba nhóm theo tốc độ suy giảm. Nhóm suy giảm chậm gồm XGBoost (slope 0,26--0,54) và E-STGCN (slope 0,21--0,68); nhóm suy giảm trung bình gồm ST-XLinear (slope 0,35--0,75); và nhóm suy giảm nhanh gồm XLinear, iTransformer và Ensemble (slope 0,37--0,84).

XGBoost duy trì slope thấp nhất trên block~5 và block~7, nhưng tại block~30 miền Nam, slope tăng lên 0,54 --- cao hơn E-STGCN (0,30) và iTransformer (0,44). Hiện tượng này cho thấy khi dữ liệu huấn luyện đủ lớn, lợi thế của hệ thống đặc trưng thủ công giảm bớt tại miền Nam --- nơi phân phối phức tạp hơn và các kiến trúc học sâu có cơ hội học được biểu diễn phù hợp.

E-STGCN ghi nhận slope thấp nhất tại block~30 (0,21 miền Bắc, 0,30 miền Nam), nhất quán với phân tích $\Delta$MAE ở Mục 3.3.1. Slope thấp phản ánh cơ chế tích chập đồ thị duy trì chất lượng dự báo ổn định ở tầm nhìn xa nhờ thông tin không gian giữa các trạm. Slope ở miền Bắc (0,21) thấp hơn miền Nam (0,30), xác nhận tương quan không gian cao hơn ở miền Bắc hỗ trợ cơ chế tích chập đồ thị hiệu quả hơn.

iTransformer có đặc trưng thú vị: slope giảm mạnh từ block~7 sang block~30 (miền Nam: từ 0,82 xuống 0,44, giảm 46\%), lớn hơn bất kỳ mô hình nào khác. Điều này nhất quán với bước nhảy $R^2$ đã phân tích ở Mục 3.5 và phản ánh cơ chế attention vượt qua ngưỡng dữ liệu tối thiểu tại block~30.

\section{Đánh giá phương pháp kết hợp mô hình}

Phương pháp Ensemble sử dụng trung bình trọng số tối ưu trên tập xác thực được kỳ vọng cải thiện hiệu suất so với từng mô hình thành phần. Hình~\ref{fig:ensemble_delta} trình bày chênh lệch MAE giữa Ensemble và mô hình thành phần có MAE thấp nhất ($\Delta_{\text{Ensemble}} = \text{MAE}_{\text{Ensemble}} - \text{MAE}_{\text{min}}$) trên toàn bộ cấu hình: giá trị dương nghĩa là Ensemble kém hơn mô hình thành phần tốt nhất.

\begin{figure}[H]
\centering
\includegraphics[width=\textwidth]{ensemble_delta.png}
\caption{Chênh lệch MAE giữa Ensemble và mô hình thành phần tốt nhất --- giá trị dương cho thấy Ensemble kém hơn}
\label{fig:ensemble_delta}
\end{figure}

Kết quả cho thấy Ensemble không cải thiện hiệu suất so với mô hình thành phần tốt nhất trên đại đa số cấu hình. $\Delta_{\text{Ensemble}}$ dương trên 28/30 cấu hình (6 block $\times$ 5 tầm nhìn), với giá trị dao động từ 0,01 đến 3,58~$\mu g/m^3$. Chỉ có 2 cấu hình ghi nhận $\Delta$ gần bằng 0, tức Ensemble ngang bằng mô hình tốt nhất. Không có cấu hình nào Ensemble vượt trội rõ rệt.

Phân tích nguyên nhân thất bại của Ensemble cần xem xét ba yếu tố. Yếu tố thứ nhất là tương quan sai số cao giữa các mô hình thành phần. Năm mô hình (XGBoost, XLinear, iTransformer, E-STGCN, ST-XLinear) đều được huấn luyện trên cùng bộ dữ liệu và cùng phân chia train/validation/test, dẫn đến các mô hình mắc lỗi tương tự trên cùng khoảng thời gian (các đợt ô nhiễm cực đoan, các thời điểm chuyển mùa). Khi tương quan sai số cao, phương pháp trung bình trọng số không thể triệt tiêu sai số bằng đa dạng hóa.

Yếu tố thứ hai là sự chênh lệch hiệu suất lớn giữa XGBoost và các mô hình còn lại. Tại block~5 và block~7, XGBoost có MAE thấp hơn rõ rệt so với bốn mô hình còn lại, đặc biệt ở tầm nhìn dài. Cơ chế tối ưu trọng số trên tập xác thực phát hiện chênh lệch này và gán trọng số cao nhất cho XGBoost (hoặc mô hình tốt nhất), khiến Ensemble suy biến thành một mô hình đơn lẻ có chất lượng bị pha loãng bởi phần đóng góp nhỏ từ các mô hình kém hơn.

Yếu tố thứ ba là cơ chế kết hợp tuyến tính không phù hợp với bài toán. Trung bình trọng số giả định sai số các mô hình thành phần xấp xỉ tuyến tính và có phương sai đồng nhất --- hai giả định không thỏa mãn khi phân phối $PM_{2.5}$ có đuôi dài. Các phương pháp kết hợp phi tuyến (stacking, meta-learning) có thể khắc phục hạn chế này nhưng không thuộc phạm vi nghiên cứu hiện tại.

\section{Chi phí tính toán}

Hiệu suất mô hình cần được đánh giá cùng với chi phí tính toán để xác định phương án tối ưu cho từng ngữ cảnh ứng dụng. Bảng~\ref{tab:training_time} trình bày thời gian huấn luyện tổng (tính bằng giây) cho toàn bộ 5 tầm nhìn $\times$ 2 vùng trên mỗi chiến lược chia dữ liệu, và Hình~\ref{fig:training_time} trực quan hóa so sánh.

\begin{table}[H]
\centering
\caption{Thời gian huấn luyện tổng (giây) --- 5 tầm nhìn $\times$ 2 vùng}
\label{tab:training_time}
\small
\begin{tabular}{l|rrr|r}
\toprule
\textbf{Mô hình} & Block 5 & Block 7 & Block 30 & Trung bình \\
\midrule
XGBoost     & \textbf{50} & \textbf{87} & \textbf{78} & \textbf{72} \\
E-STGCN     & 173 & 197 & 180 & 183 \\
Ensemble    & 1.956 & 1.446 & 1.694 & 1.699 \\
iTransformer & 1.746 & 1.879 & 2.013 & 1.879 \\
ST-XLinear  & 2.285 & 2.817 & 2.616 & 2.573 \\
XLinear     & 3.888 & 3.626 & 3.210 & 3.575 \\
\bottomrule
\end{tabular}
\end{table}

\begin{figure}[H]
\centering
\includegraphics[width=\textwidth]{training_time.png}
\caption{Thời gian huấn luyện theo mô hình và chiến lược chia dữ liệu}
\label{fig:training_time}
\end{figure}

Sáu mô hình phân thành ba nhóm chi phí rõ rệt. Nhóm chi phí thấp gồm XGBoost (50--87 giây) và E-STGCN (173--197 giây), nhanh hơn nhóm còn lại 10--40 lần. Nhóm chi phí trung bình gồm Ensemble (1.446--1.956 giây) và iTransformer (1.746--2.013 giây). Nhóm chi phí cao gồm ST-XLinear (2.285--2.817 giây) và XLinear (3.210--3.888 giây).

XGBoost nhanh nhất nhờ cơ chế gradient boosting trên cây quyết định, vốn tận dụng hiệu quả phần cứng CPU và không yêu cầu GPU. Trung bình 72 giây cho 10 mô hình (5 tầm nhìn $\times$ 2 vùng) tương đương khoảng 7 giây mỗi mô hình --- cho phép tái huấn luyện hàng giờ trong hệ thống thời gian thực. E-STGCN nhanh thứ hai (trung bình 183 giây) nhờ kiến trúc nhỏ gọn với ít tham số so với iTransformer; thời gian này vẫn phù hợp cho cập nhật hàng ngày.

XLinear có thời gian huấn luyện cao nhất (trung bình 3.575 giây, tức gần 60 phút), mặc dù kiến trúc tuyến tính đơn giản. Nguyên nhân là XLinear huấn luyện riêng biệt cho 12 trạm trong mỗi vùng (tổng 12 $\times$ 5 $\times$ 2 = 120 mô hình con), và mỗi mô hình con cần tìm kiếm siêu tham số trên tập xác thực. ST-XLinear kế thừa thời gian cao này do chia sẻ cùng quy trình huấn luyện nền XLinear, cộng thêm thành phần không gian.

Khi kết hợp hiệu suất và chi phí, ba kịch bản triển khai có thể được xem xét. Kịch bản thứ nhất là hệ thống thời gian thực cần cập nhật liên tục (dưới 5 phút): XGBoost là phương án phù hợp duy nhất, với thời gian huấn luyện dưới 90 giây và hiệu suất dẫn đầu hoặc gần dẫn đầu trên mọi cấu hình. Kịch bản thứ hai là hệ thống cập nhật hàng ngày (dưới 30 phút): E-STGCN trở nên cạnh tranh với chi phí khoảng 3 phút và hiệu suất tại tầm nhìn dài cạnh tranh với XGBoost, đặc biệt trên block~30. Kịch bản thứ ba là nghiên cứu và đánh giá (không giới hạn thời gian): iTransformer đáng xem xét nhờ mức cải thiện lớn tại block~30 mặc dù chi phí khoảng 30 phút, nhưng XLinear và ST-XLinear khó biện minh do chi phí cao nhưng hiệu suất không vượt trội rõ rệt so với các phương án rẻ hơn.

\section{Ứng dụng web dự báo chất lượng không khí}

Mục tiêu thứ năm của đồ án là triển khai ứng dụng web cho phép người dùng tra cứu dự báo $PM_{2.5}$ trực tuyến. Ứng dụng được xây dựng trên nền tảng Streamlit (Python) và tích hợp mô hình XGBoost block~7 để thực hiện dự báo thời gian thực tại 12 trạm quan trắc.

\subsection{Kiến trúc ứng dụng}

Ứng dụng hoạt động theo luồng xử lý ba bước. Bước thứ nhất, khi người dùng chọn trạm quan trắc, hệ thống xác định trạm thuộc cụm miền Bắc hay miền Nam dựa trên bảng ánh xạ cố định (6 trạm Bắc: ID 1, 4, 5, 16, 17, 27; 6 trạm Nam: ID 7, 18, 24, 30, 31, 32). Bước thứ hai, module suy luận (inference) đọc toàn bộ dữ liệu lịch sử của trạm, xây dựng 134 đặc trưng đầu vào (bao gồm độ trễ, trung bình trượt, đặc trưng thời gian, đặc trưng hóa học khí quyển và dữ liệu thời tiết tương lai), sau đó nạp vào mô hình XGBoost tương ứng với cụm và tầm nhìn dự báo. Bước thứ ba, giá trị dự báo được chuyển ngược (inverse-transform) về đơn vị $\mu g/m^3$ bằng bộ chuẩn hóa đã lưu trong quá trình huấn luyện, và hiển thị lên giao diện.

Mô hình XGBoost block~7 được chọn cho ứng dụng vì ba lý do: thời gian suy luận nhanh (dưới 1 giây cho 5 tầm nhìn), hiệu suất dẫn đầu hoặc gần dẫn đầu trên đa số cấu hình (như đã phân tích ở các mục trước), và không yêu cầu GPU để chạy suy luận.

\subsection{Giao diện tổng quan}

Hình~\ref{fig:header_web} trình bày giao diện tổng quan của ứng dụng.

\begin{figure}[H]
\centering
\includegraphics[width=\textwidth]{Header_web.png}
\caption{Giao diện tổng quan --- thanh điều khiển bên trái, thẻ chỉ số và biểu đồ lịch sử 24 giờ}
\label{fig:header_web}
\end{figure}

Giao diện chia thành hai phần chính. Phần bên trái là thanh điều khiển cho phép chọn trạm quan trắc từ danh sách 12 trạm, mỗi trạm hiển thị kèm nhãn cụm (Bắc/Nam), tên quận huyện và tỉnh thành. Khi chọn trạm, hệ thống hiển thị tọa độ GPS và loại mô hình XGBoost tương ứng (North hoặc South). Phần bên phải hiển thị năm thẻ chỉ số tại thời điểm mới nhất: $PM_{2.5}$, $PM_{10}$, nhiệt độ, độ ẩm và tốc độ gió. Bên dưới là nhãn trạng thái chất lượng không khí với mã màu theo thang AQI (xanh lá = Tốt, vàng = Trung bình, cam = Không tốt cho nhóm nhạy cảm, đỏ = Không tốt). Biểu đồ chính hiển thị diễn biến $PM_{2.5}$ trong 24 giờ gần nhất dưới dạng đường cong với ba đường ngưỡng AQI tham chiếu, kết hợp biểu đồ nhiệt độ ở phần dưới để người dùng nhận biết mối tương quan giữa hai yếu tố.

\subsection{Hiển thị các chỉ số ô nhiễm}

Hình~\ref{fig:chiso_24} trình bày khu vực hiển thị xu hướng 24 giờ của sáu chỉ số ô nhiễm.

\begin{figure}[H]
\centering
\includegraphics[width=\textwidth]{Chi__so_24.png}
\caption{Biểu đồ thu nhỏ xu hướng 24 giờ của sáu chỉ số ô nhiễm}
\label{fig:chiso_24}
\end{figure}

Mỗi chỉ số ($PM_{2.5}$, $PM_{10}$, CO, $O_3$, $NO_2$, $SO_2$) được hiển thị trong một ô riêng biệt gồm ba thành phần: giá trị hiện tại ở dạng số lớn, mũi tên chỉ hướng thay đổi so với giờ trước (tam giác lên màu cam = tăng, tam giác xuống màu xanh = giảm), và biểu đồ thu nhỏ dạng đường cong thể hiện xu hướng 24 giờ. Thiết kế này cho phép người dùng nhanh chóng nhận biết trạng thái hiện tại và chiều hướng biến động của từng chỉ số mà không cần phân tích biểu đồ chi tiết.

\subsection{Khu vực dự báo}

Hình~\ref{fig:prediction_xgboost} trình bày khu vực dự báo $PM_{2.5}$ của ứng dụng.

\begin{figure}[H]
\centering
\includegraphics[width=\textwidth]{Prediction_XGBoost.png}
\caption{Khu vực dự báo $PM_{2.5}$ --- năm thẻ dự báo và biểu đồ nối lịch sử--dự báo}
\label{fig:prediction_xgboost}
\end{figure}

Khu vực dự báo gồm hai thành phần. Thành phần thứ nhất là năm thẻ hiển thị giá trị $PM_{2.5}$ dự báo tại các mốc $T{+}1$, $T{+}3$, $T{+}6$, $T{+}12$ và $T{+}24$ giờ. Mỗi thẻ hiển thị giá trị dự báo ($\mu g/m^3$), nhãn mức chất lượng không khí và loại dự báo (ngắn hạn, trung ngắn, trung hạn, dài hạn). Màu sắc giá trị và viền thẻ thay đổi động theo thang AQI, giúp người dùng nhận biết ngay mức độ ô nhiễm dự kiến.

Thành phần thứ hai là biểu đồ kết hợp lịch sử và dự báo. Đường liền nét thể hiện giá trị $PM_{2.5}$ thực tế trong 24 giờ gần nhất, đường nét đứt thể hiện giá trị dự báo tại năm mốc thời gian tương lai. Đường dọc phân chia ranh giới giữa dữ liệu lịch sử và dự báo, kết hợp ba đường ngưỡng AQI nằm ngang để người dùng đối chiếu giá trị dự báo với các mức cảnh báo. Mỗi điểm dự báo trên biểu đồ có màu sắc tương ứng với mức AQI, tạo sự nhất quán với các thẻ dự báo phía trên.

Bên dưới khu vực dự báo, ứng dụng hiển thị bảng đối chiếu sáu mức AQI theo $PM_{2.5}$ (từ Tốt đến Nguy hiểm) với khoảng giá trị tương ứng, giúp người dùng hiểu ý nghĩa các mã màu được sử dụng xuyên suốt giao diện.

\section{Thảo luận}

Chương này đã trình bày phân tích toàn diện kết quả thí nghiệm của sáu mô hình dự báo $PM_{2.5}$ trên 180 cấu hình (6 mô hình $\times$ 5 tầm nhìn $\times$ 3 block $\times$ 2 vùng). Phần thảo luận đánh giá mức độ đạt được từng mục tiêu nghiên cứu đã đề ra trong lời mở đầu, đồng thời phân tích các ưu điểm và hạn chế của nghiên cứu.

\subsection{Đánh giá mức độ đạt mục tiêu}

Mục tiêu thứ nhất --- xây dựng pipeline thu thập và tiền xử lý dữ liệu --- đã được hoàn thành. Từ 32 trạm quan trắc ban đầu, quy trình sàng lọc chất lượng dữ liệu đã chọn được 12 trạm (6 miền Bắc, 6 miền Nam) với dữ liệu theo giờ giai đoạn 01/2023 -- 12/2025, đảm bảo bao phủ đầy đủ các chu kỳ mùa. Pipeline tiền xử lý bao gồm phát hiện lỗi cảm biến đóng băng, xử lý ngoại lai bằng IQR, nội suy giá trị thiếu và chuẩn hóa theo 8 chiến lược khác nhau tùy nhóm biến. Kết quả benchmark cho thấy dữ liệu sau tiền xử lý đủ chất lượng để huấn luyện cả sáu kiến trúc mô hình trên mọi cấu hình thí nghiệm.

Mục tiêu thứ hai --- thiết kế đặc trưng kỹ thuật dựa trên kiến thức hóa học khí quyển --- đã được triển khai. Hệ thống 134 đặc trưng đầu vào bao gồm 6 đặc trưng mới phản ánh khả năng oxy hóa, nguy cơ hình thành aerosol sulfate và mức độ ổn định nhiệt khí quyển, kết hợp với đặc trưng trễ, thống kê trượt và mã hóa chu kỳ. Hiệu suất XGBoost --- mô hình sử dụng trực tiếp toàn bộ hệ thống đặc trưng --- dẫn đầu trên block~5 và block~7, cho thấy hệ thống đặc trưng thủ công cung cấp thông tin dự báo có giá trị, đặc biệt khi dữ liệu huấn luyện hạn chế.

Mục tiêu thứ ba --- huấn luyện và so sánh sáu mô hình --- là nội dung chính của chương này. Phân tích cho thấy không có mô hình nào vượt trội tuyệt đối trên mọi cấu hình. XGBoost đạt hiệu suất cao nhất tại block~5 và block~7 với MAE thấp nhất, slope suy giảm chậm nhất và thời gian huấn luyện nhanh nhất (trung bình 72 giây). Tại block~30, các mô hình học sâu thu hẹp khoảng cách: iTransformer cải thiện $R^2$ lên 0,34 so với block~7, E-STGCN đạt slope thấp nhất (0,21 miền Bắc) nhờ tích chập đồ thị. Kết quả này phản ánh sự đánh đổi giữa yêu cầu dữ liệu và khả năng học biểu diễn: mô hình thống kê cổ điển cần ít dữ liệu nhưng giới hạn ở đặc trưng thủ công, trong khi mô hình học sâu cần nhiều dữ liệu nhưng có tiềm năng học được biểu diễn phức tạp hơn.

Mục tiêu thứ tư --- đề xuất cơ chế Dynamic Graph --- đã được triển khai trong mô hình ST-XLinear. Tuy nhiên, kết quả cho thấy ST-XLinear không vượt trội so với XLinear (mô hình nền không có thành phần không gian) trên đa số cấu hình: $R^2$ tại block~30 miền Bắc $T{+}24$ của ST-XLinear (0,37) thấp hơn XLinear (0,40). Nguyên nhân có thể bao gồm: mạng lưới 6 trạm trong mỗi cụm quá thưa để ma trận kề động tạo ra giá trị bổ sung đáng kể, và cơ chế cập nhật theo hướng gió có thể gây nhiễu khi hướng gió thay đổi nhanh trong ngày. E-STGCN sử dụng ma trận kề tĩnh (dựa trên khoảng cách địa lý) lại cho kết quả tốt hơn tại block~30, gợi ý rằng quan hệ không gian ổn định có thể phù hợp hơn cho mạng lưới thưa.

Mục tiêu thứ năm --- triển khai ứng dụng web dự báo --- đã hoàn thành như trình bày ở Mục 3.9. Ứng dụng cho phép tra cứu dự báo $PM_{2.5}$ tại 12 trạm với năm mốc thời gian, tích hợp mô hình XGBoost block~7 với thời gian suy luận dưới 1 giây.

\subsection{Ưu điểm của nghiên cứu}

Nghiên cứu có ba ưu điểm chính. Thứ nhất, quy mô benchmark toàn diện với 180 cấu hình thí nghiệm cho phép đánh giá hiệu suất mô hình trên nhiều chiều: tầm nhìn dự báo, chiến lược chia dữ liệu và vùng địa lý. Cách tiếp cận này giảm nguy cơ kết luận thiên lệch do chỉ đánh giá trên một cấu hình duy nhất --- một hạn chế phổ biến trong các nghiên cứu dự báo chất lượng không khí.

Thứ hai, phân tích đa chiều kết hợp cả bốn chỉ số (MAE, RMSE, $R^2$, MAPE) cùng các phân tích bổ sung (slope suy giảm, biểu đồ hộp, biểu đồ radar) cho phép nhận diện các đặc trưng mà một chỉ số đơn lẻ không thể phát hiện. Ví dụ, E-STGCN có $R^2$ cao nhưng MAE cũng cao --- mâu thuẫn này chỉ được giải thích khi phân tích kết hợp nhiều chỉ số.

Thứ ba, nghiên cứu cung cấp phân tích nguyên nhân cho các quan sát thay vì chỉ báo cáo số liệu. Sự khác biệt hiệu suất giữa hai vùng được giải thích bằng đặc trưng phân phối dữ liệu, điều kiện khí tượng và mật độ không gian trạm; ảnh hưởng của block split được giải thích bằng yêu cầu dữ liệu tối thiểu của từng kiến trúc.

\subsection{Hạn chế của nghiên cứu}

Nghiên cứu tồn tại bốn hạn chế cần lưu ý. Thứ nhất, mạng lưới 12 trạm quan trắc (6 trạm mỗi cụm) là quy mô nhỏ so với diện tích lãnh thổ. Mật độ trạm thưa khiến các mô hình không gian (E-STGCN, ST-XLinear) không phát huy được tối đa tiềm năng kiến trúc, và kết quả chưa chắc tổng quát hóa được cho mạng lưới dày hơn.

Thứ hai, phương pháp Ensemble sử dụng trung bình trọng số tuyến tính đã cho thấy hiệu quả hạn chế. Các phương pháp kết hợp phi tuyến như stacking (sử dụng mô hình bậc hai để học cách kết hợp đầu ra của các mô hình thành phần) hoặc phương pháp kết hợp theo miền (gán trọng số khác nhau cho từng khoảng giá trị $PM_{2.5}$) có thể cải thiện kết quả nhưng chưa được khảo sát.

Thứ ba, đánh giá hiện tại sử dụng phương pháp block-based split thay vì phương pháp xác thực chéo chuỗi thời gian (time series cross-validation) với nhiều lần lặp. Mặc dù ba cấu hình block (5, 7, 30 ngày) cung cấp một mức độ đánh giá tính ổn định, kết quả có thể bị ảnh hưởng bởi đặc điểm cụ thể của các giai đoạn thời gian được chọn làm tập kiểm thử.

Thứ tư, nghiên cứu chưa so sánh với các phương pháp dự báo dựa trên mô phỏng vật lý (CMAQ, WRF-Chem) hoặc các mô hình ngôn ngữ lớn được tinh chỉnh cho chuỗi thời gian. Việc bổ sung các đường cơ sở này sẽ giúp định vị rõ hơn vị trí của các phương pháp học máy trong bài toán dự báo chất lượng không khí tại Việt Nam.
